\fchapter{Ej 9 Pág 80}

\fsection{\boldit{\textcolor{waveBlue2}{Actividad.}}} 
\subsection*{Reflexiona sobre los beneficios que obtienen algunos colectivos vulnerables, como los discapacitados, en el uso de la tecnología y cómo se contribuye a la reducción de las desigualdades que sufren en su día a día.}

\begin{solution}
La tecnología es una herramienta clave para reducir las desigualdades. En el caso de las personas con discapacidad, su impacto es transformador:
\begin{enumerate}
	\item \textbf{Autonomía y vida independiente:}
		\begin{itemize}
			\item{Ayudas técnicas: Sillas de ruedas eléctricas, prótesis robóticas y domótica (control por voz de la casa) permiten un mayor control del entorno.}
			\item{Apps de ayuda: Aplicaciones para la comunicación (texto a voz) y la navegación (mapas accesibles) facilitan la vida diaria.}
		\end{itemize}
	\item \textbf{Acceso a la información, educación y comunicación:
}
		\begin{itemize}
			\item{Lectores de pantalla y magnificadores para personas con discapacidad visual.}
			\item{Subtítulos automáticos y intérpretes de lengua de signos en línea para personas con discapacidad auditiva.}
			\item{Plataformas educativas accesibles que rompen barreras para el aprendizaje.}
		\end{itemize}

	\item \textbf{Inclusión social y laboral:}
		\begin{itemize}
			\item{Herramientas de teletrabajo y comunicación remota que abren nuevas oportunidades laborales.}
			\item{Redes sociales y foros que permiten conectar con otras personas, reduciendo el aislamiento.}
		\end{itemize}

\end{enumerate}


\end{solution}
